\documentclass[aspectratio=169,11pt]{beamer}
\usetheme{default}
\setbeamertemplate{navigation symbols}{}
\setbeamertemplate{footline}{%
  \hfill\insertframenumber/\inserttotalframenumber\hspace*{4pt}\vspace{2pt}%
}
\usepackage{lmodern}
\usepackage[most]{tcolorbox}
\usepackage{listings}
\usepackage{tikz}
\usetikzlibrary{arrows.meta,positioning,shapes.geometric,fit,backgrounds}
\usepackage{booktabs}
\usepackage{hyperref}

\definecolor{lcgreen}{HTML}{2E7D32}
\definecolor{lcblue}{HTML}{1565C0}
\definecolor{lcorange}{HTML}{EF6C00}
\definecolor{lcred}{HTML}{C62828}
\definecolor{codebg}{HTML}{F5F5F5}
\definecolor{docgray}{gray}{0.55}

\newtcolorbox{conceptbox}[1][]{colback=yellow!20,colframe=yellow!60!black,title=#1,fonttitle=\bfseries,top=2pt,bottom=2pt,left=4pt,right=4pt}
\newtcolorbox{warningbox}[1][]{colback=red!15,colframe=red!60!black,title=#1,fonttitle=\bfseries,top=2pt,bottom=2pt,left=4pt,right=4pt}
\newtcolorbox{tipbox}[1][]{colback=green!10,colframe=green!50!black,title=#1,fonttitle=\bfseries,top=2pt,bottom=2pt,left=4pt,right=4pt}
\newtcolorbox{codebox}[1][]{colback=blue!10,colframe=blue!40!black,title=#1,fonttitle=\bfseries,top=2pt,bottom=2pt,left=4pt,right=4pt}

\lstset{
  basicstyle=\ttfamily\scriptsize,
  backgroundcolor=\color{codebg},
  breaklines=true,
  frame=single,
  rulecolor=\color{blue!40!black},
  keywordstyle=\color{lcblue}\bfseries,
  stringstyle=\color{lcgreen},
  commentstyle=\color{docgray}\itshape,
  language=Python,
  showstringspaces=false,
  tabsize=2,
  xleftmargin=2pt,
  xrightmargin=2pt,
  aboveskip=2pt,
  belowskip=2pt,
}

\newcommand{\docref}[1]{{\scriptsize\color{docgray}[Docs: #1]}}

\title{LangGraph v1.0.x}
\subtitle{Section 03: Conditional Routing \& Branching}
\author{Agentic AI Foundations}
\date{March 2026}

\begin{document}

\begin{frame}
\titlepage
\end{frame}

\section{Mental Model}
\begin{frame}[t]{Making Graphs Dynamic with Conditional Logic}
\setlength{\itemsep}{1pt}
\vspace{0.1cm}

\begin{conceptbox}[Conditional Routing = Agentic Flow]
Simple graphs are linear. Agents need to make \textit{decisions}. Conditional edges route based on state.
\end{conceptbox}

\vspace{0.15cm}

\begin{tipbox}[Examples]
\setlength{\itemsep}{1pt}
\begin{itemize}
  \item Tool use agent: ``Did LLM call a tool? Route to tool executor.''
  \item Classifier: ``Route to approve or reject path.''
  \item Retry loop: ``Did it fail? Go back to try again.''
\end{itemize}
\end{tipbox}
\end{frame}

\section{Why Conditional Routing}
\begin{frame}[t]{Why Conditional Routing: Decision-Making}
\setlength{\itemsep}{1pt}
\vspace{0.1cm}

\begin{warningbox}[Linear Isn't Enough]
Linear graphs execute the same sequence every time. Agents need to \textit{decide} where to go next based on state.
\end{warningbox}

\vspace{0.15cm}

\begin{conceptbox}[Conditional Edges Enable]
\setlength{\itemsep}{1pt}
\begin{itemize}
  \item Branching paths (A -> B or C or D)
  \item Looping (A -> A again)
  \item Early exit (-> END)
  \item Dynamic workflows
\end{itemize}
\end{conceptbox}
\end{frame}

\section{Router Functions}
\begin{frame}[fragile,t]{Router Functions: Read State, Return Next Node}
\vspace{0.1cm}

\begin{conceptbox}[Router Signature]
A router function reads state and returns a \textit{string} naming the next node.
\end{conceptbox}

\vspace{0.1cm}

\begin{codebox}[Simple Router]
\begin{lstlisting}
def my_router(state: State) -> str:
    if state["score"] > 0.5:
        return "approve"
    else:
        return "reject"
\end{lstlisting}
\end{codebox}

\vspace{0.1cm}

{\footnotesize Router returns a string. The graph looks it up in a path map.}
\end{frame}

\section{add_conditional_edges}
\begin{frame}[fragile,t]{add\_conditional\_edges: Source, Router, Path Map}
\vspace{0.05cm}

\begin{codebox}[Conditional Edges]
\begin{lstlisting}
graph.add_conditional_edges(
    "source_node",          # From
    my_router,              # Router function
    {
        "approve": "approve_node",
        "reject": "reject_node"
    }
)
\end{lstlisting}
\end{codebox}

\vspace{0.1cm}

\begin{tipbox}[Three Parts]
\setlength{\itemsep}{1pt}
\begin{itemize}
  \item Source node (where we're coming from)
  \item Router function (logic to decide)
  \item Path map (string -> node name)
\end{itemize}
\end{tipbox}
\end{frame}

\section{Multiple Branches}
\begin{frame}[fragile,t]{Multiple Conditional Branches}
\vspace{0.05cm}

\begin{codebox}[Router with Many Paths]
\begin{lstlisting}
def route_action(state: State) -> str:
    action = state["action"]
    if action == "approve":
        return "approve"
    elif action == "reject":
        return "reject"
    else:
        return "review"

graph.add_conditional_edges(
    "classifier",
    route_action,
    {
        "approve": "approve_node",
        "reject": "reject_node",
        "review": "review_node"
    }
)
\end{lstlisting}
\end{codebox}
\end{frame}

\section{Default Routing}
\begin{frame}[fragile,t]{Default/Fallback Routing}
\vspace{0.1cm}

\begin{conceptbox}[When Router Returns Unknown Value]
Set a default fallback path.
\end{conceptbox}

\vspace{0.1cm}

\begin{codebox}[With Default]
\begin{lstlisting}
graph.add_conditional_edges(
    "router",
    router_fn,
    {
        "path_a": "node_a",
        "path_b": "node_b"
    },
    default="node_default"  # Fallback
)
\end{lstlisting}
\end{codebox}

\vspace{0.15cm}

{\footnotesize If router returns something not in the map, default catches it.}
\end{frame}

\section{Cycles and Loops}
\begin{frame}[t]{Cycles and Loops: Agentic Feedback}
\setlength{\itemsep}{1pt}
\vspace{0.1cm}

\begin{conceptbox}[Loop Pattern]
A node can route back to itself (or an earlier node) for iteration.
\end{conceptbox}

\vspace{0.15cm}

\begin{tipbox}[Use Cases]
\setlength{\itemsep}{1pt}
\begin{itemize}
  \item Agent think->act→observe→think cycle
  \item Retry logic (failed? try again)
  \item Refinement loops
\end{itemize}
\end{tipbox}

\vspace{0.15cm}

{\footnotesize LangGraph handles cycles natively. No special syntax needed.}
\end{frame}

\begin{frame}[fragile,t]{Breaking Out of Loops: Route to END}
\vspace{0.1cm}

\begin{conceptbox}[Termination Condition]
A loop must have a condition that routes to END, or it runs forever!
\end{conceptbox}

\vspace{0.1cm}

\begin{codebox}[Loop with Exit]
\begin{lstlisting}
def should_continue(state: State) -> str:
    if state["done"]:
        return "done"  # Route to END
    else:
        return "loop"  # Route back

graph.add_conditional_edges(
    "agent",
    should_continue,
    {"done": END, "loop": "agent"}
)
\end{lstlisting}
\end{codebox}
\end{frame}

\section{tools_condition Helper}
\begin{frame}[fragile,t]{tools\_condition: Routing on Tool Use}
\vspace{0.05cm}

\begin{conceptbox}[Common Pattern]
Did the LLM call tools? Route to tool executor, else to END.
\end{conceptbox}

\vspace{0.1cm}

\begin{codebox}[Using tools\_condition]
\begin{lstlisting}
from langgraph.prebuilt import tools_condition

graph.add_conditional_edges(
    "agent",
    tools_condition,
    {
        "tools": "execute_tools",
        "end": END
    }
)
\end{lstlisting}
\end{codebox}

\vspace{0.15cm}

{\footnotesize \texttt{tools\_condition} checks if last message has tool\_calls.}
\end{frame}

\section{Agent Loop Diagram}
\begin{frame}[t]{Complete Agent Loop: Model -> Tools → Model → END}
\vspace{0.1cm}

\begin{tikzpicture}[node distance=1.2cm, scale=0.95, every node/.style={font=\scriptsize}]
  \node[draw, fill=lcblue!30, rounded rectangle] (s) {START};
  \node[draw, fill=lcgreen!30, rounded rectangle, right=of s] (llm) {LLM};
  \node[draw, fill=lcgreen!30, rounded rectangle, below=of llm] (tool) {Tool};
  \node[draw, fill=lcred!30, rounded rectangle, right=of llm] (end) {END};

  \path[->, thick] (s) -- (llm);
  \path[->, thick] (llm) -- node[right, font=\tiny] {calls tools?} (tool);
  \path[->, thick] (tool) -- (llm);
  \path[->, thick] (llm) -- node[above, font=\tiny] {no tools} (end);
\end{tikzpicture}

\vspace{0.15cm}

The agent loop: think (LLM) -> decide (router) → act (tools) → think again.
\end{frame}

\section{Branching Patterns}
\begin{frame}[t]{Branching Patterns: Fan-Out and Fan-In}
\setlength{\itemsep}{1pt}
\vspace{0.1cm}

\begin{tipbox}[Fan-Out]
One node routes to multiple nodes in parallel (or sequence).
\end{tipbox}

\vspace{0.1cm}

\begin{tipbox}[Fan-In]
Multiple nodes converge to one. All must complete before proceeding.
\end{tipbox}

\vspace{0.15cm}

\begin{conceptbox}[Use Cases]
\setlength{\itemsep}{1pt}
\begin{itemize}
  \item Fan-out: process multiple paths
  \item Fan-in: aggregate results
\end{itemize}
\end{conceptbox}
\end{frame}

\section{Map-Reduce Pattern}
\begin{frame}[fragile,t]{Map-Reduce: Process Items, Then Aggregate}
\vspace{0.05cm}

\begin{conceptbox}[Pattern]
Fan-out to process each item, then fan-in to aggregate results.
\end{conceptbox}

\vspace{0.1cm}

\begin{codebox}[Idea]
\begin{lstlisting}
# State: items: [1, 2, 3, ...]
# Map: process each
# Reduce: combine results
\end{lstlisting}
\end{codebox}

\vspace{0.15cm}

{\footnotesize LangGraph supports this with conditional edges + loops.}
\end{frame}

\section{Send API}
\begin{frame}[fragile,t]{Send API for Dynamic Fan-Out}
\vspace{0.1cm}

\begin{conceptbox}[Dynamic Routing]
For \textit{unknown} number of paths, use Send() to create edges dynamically.
\end{conceptbox}

\vspace{0.1cm}

\begin{codebox}[Send Example]
\begin{lstlisting}
from langgraph.graph import Send

def route_to_many(state):
    return [
        Send("worker", {"item": x})
        for x in state["items"]
    ]
\end{lstlisting}
\end{codebox}

\vspace{0.15cm}

{\footnotesize Advanced pattern. Useful for variable-width fan-out.}
\end{frame}

\section{Common Mistakes}
\begin{frame}[fragile,t]{Common Mistakes: Infinite Loops}
\vspace{0.1cm}

\begin{warningbox}[Mistake 1: Infinite Loop]
\begin{lstlisting}
# WRONG: routes back to self forever
graph.add_conditional_edges(
    "agent",
    my_router,
    {"continue": "agent"}
    # Missing exit condition
)
\end{lstlisting}
\end{warningbox}

\vspace{0.15cm}

\begin{tipbox}[Fix]
Always have a condition that routes to END!
\end{tipbox}
\end{frame}

\begin{frame}[fragile,t]{Common Mistakes: Missing Paths in Map}
\vspace{0.1cm}

\begin{warningbox}[Mistake 2: Router Returns Unregistered Path]
\begin{lstlisting}
# Router returns "unknown"
def router(state): return "unknown"

# But path map doesn't have it
graph.add_conditional_edges(
    "node",
    router,
    {"path_a": "a", "path_b": "b"}
    # "unknown" not mapped!
)
\end{lstlisting}
\end{warningbox}

\vspace{0.15cm}

\begin{tipbox}[Fix]
Use default, or ensure router only returns known paths.
\end{tipbox}
\end{frame}

\section{Advanced Patterns}
\begin{frame}[t]{Advanced Patterns Summary}
\setlength{\itemsep}{1pt}
\vspace{0.1cm}

\begin{tipbox}[Sequential Branching]
Node A -> (B or C or D) in sequence.
\end{tipbox}

\vspace{0.1cm}

\begin{tipbox}[Parallel Branching]
Node A -> multiple nodes at once (with Send).
\end{tipbox}

\vspace{0.1cm}

\begin{tipbox}[Looping]
Node A -> (A again or continue).
\end{tipbox}

\vspace{0.1cm}

\begin{tipbox}[Hierarchical]
Subgraphs within graphs (advanced).
\end{tipbox}
\end{frame}

\section{Summary}
\begin{frame}[t,shrink=3]{Summary: Conditional Routing \& Branching}
\setlength{\itemsep}{1pt}
\vspace{0.05cm}

\begin{conceptbox}[Core Concepts]
\setlength{\itemsep}{0.5pt}
\begin{itemize}
  \item Router functions: read state, return node name
  \item add\_conditional\_edges: source, router, path map
  \item Cycles: loops for agent reasoning
  \item Fan-out/fan-in: branching and aggregation
  \item Send API: dynamic routing
\end{itemize}
\end{conceptbox}

\vspace{0.1cm}

\begin{warningbox}[Critical]
\textit{Always} have an exit condition for loops!
\end{warningbox}

\vspace{0.1cm}

\begin{tipbox}[Pattern]
Agent loop = LLM -> router → tool → loop back (or END)
\end{tipbox}
\end{frame}

\begin{frame}[t,shrink=5]{Self-Check Questions}
\setlength{\itemsep}{0pt}
\vspace{0.05cm}

{\footnotesize
\begin{enumerate}
\setlength{\itemsep}{0.5pt}
  \item What does a router function return?
  \item How do conditional edges work (3 parts)?
  \item What's the difference between fan-out and fan-in?
  \item Why do loops need an exit condition?
  \item When would you use tools\_condition?
  \item What does Send API do?
  \item How do you implement an agent loop?
\end{enumerate}

\vspace{0.05cm}

\textbf{Quick Answers:} (1) String (node name). (2) Source, router fn, path map. (3) Fan-out: one->many; fan-in: many→one. (4) Without exit, infinite loop. (5) Route on LLM tool calls. (6) Dynamic fan-out. (7) LLM→router→tool→loop/END.
}
\end{frame}

\begin{frame}[t]{Next Steps Beyond Section 03}
\setlength{\itemsep}{1pt}
\vspace{0.1cm}

\begin{tipbox}[You Now Know]
\setlength{\itemsep}{1pt}
\begin{itemize}
  \item State, nodes, edges (Section 01)
  \item Building graphs (Section 02)
  \item Conditional routing \& branching (Section 03)
\end{itemize}
\end{tipbox}

\vspace{0.15cm}

\begin{conceptbox}[Future Topics]
Persistence, tools, streaming, human-in-the-loop, deployment, evaluation.
\end{conceptbox}

\vspace{0.15cm}

You're ready to build production agentic workflows!
\end{frame}

\end{document}
